\documentclass[../../../include/open-logic-section]{subfiles}

\begin{document}

\olfileid[fr]{sth}{z}{power}
\olsection{Ensembles des parties}

Passons à un autre axiome :

\begin{axiom}[Ensemble des parties]
Pour tout ensemble $A$, l'ensemble
$\Pow{A} = \Setabs{x}{x \subseteq A}$ existe.
\[
	\forall A \exists P \forall x(x \in P \liff (\forall z \in x)z \in A)
\]
\end{axiom}

La justification est assez directe. Supposons que $A$ soit formé à
l'étape $S$. Tous les !!{element}s de $A$ étaient alors disponibles
avant $S$ (par \stagesacc). Ainsi, en raisonnant comme pour justifier
la séparation, tout sous-ensemble de $A$ est formé au plus tard à
l'étape $S$. Tous ces sous-ensembles sont donc disponibles pour être
réunis en un ensemble à n'importe quelle étape postérieure à $S$.
Or nous savons qu'une telle étape existe, puisque $S$ n'est pas la
dernière (par \stagessucc). Ainsi, $\Pow{A}$ existe.

Voici une conséquence intéressante de l'axiome de l'ensemble des
parties :

\begin{prop}\label{thm:Products}
Quels que soient les ensembles $A, B$, leur produit cartésien
$A \times B$ existe.
\end{prop}

\begin{proof}
L'ensemble $\Pow{\Pow{A \cup B}}$ existe par l'axiome de l'ensemble
des parties et par la
\olref[sth][z][pairs]{prop:pairsconsequences}. Par séparation,
l'ensemble suivant existe donc :
\[
	C = \Setabs{z \in \Pow{\Pow{A \cup B}}}{(\exists x \in A)(\exists y \in B) z = \tuple{x, y}}.
\]
Or, pour tous $x \in A$ et $y \in B$, l'ensemble $\tuple{x, y}$ existe
par la \olref[sth][z][pairs]{prop:pairsconsequences}. De plus, puisque
$x, y \in A \cup B$, nous avons
$\{x\}, \{x, y\} \in \Pow{A \cup B}$, puis
$\tuple{x,y} \in \Pow{\Pow{A \cup B}}$. Ainsi, $A \times B = C$.
\end{proof}

Dans cette démonstration, l'axiome de l'ensemble des parties intervient
avec la séparation. Rien d'étonnant à cela : sans la séparation,
l'axiome de l'ensemble des parties ne serait pas très \emph{puissant}.
Après tout, c'est la séparation qui nous dit quels sous-ensembles d'un
ensemble existent et qui détermine ainsi à quel point chaque ensemble
des parties est «~gros~».

\begin{prob}
Montrez que, pour tous ensembles $A, B$ : (i) l'ensemble de toutes
les relations de domaine $A$ et d'image $B$ existe ; et (ii)
l'ensemble de toutes les fonctions de $A$ vers $B$ existe.
\end{prob}

\begin{prob}
Soit $A$ un ensemble, et soit $\sim$ une relation d'équivalence sur
$A$. Démontrez que l'ensemble des classes d'équivalence modulo
$\sim$ sur $A$, c'est-à-dire $\equivclass{A}{\sim}$, existe.
\end{prob}

\end{document}
